% Options for packages loaded elsewhere
\PassOptionsToPackage{unicode}{hyperref}
\PassOptionsToPackage{hyphens}{url}
\PassOptionsToPackage{dvipsnames,svgnames,x11names}{xcolor}
%
\documentclass[
  11pt,
]{article}
\usepackage{amsmath,amssymb}
\usepackage{iftex}
\ifPDFTeX
  \usepackage[T1]{fontenc}
  \usepackage[utf8]{inputenc}
  \usepackage{textcomp} % provide euro and other symbols
\else % if luatex or xetex
  \usepackage{unicode-math} % this also loads fontspec
  \defaultfontfeatures{Scale=MatchLowercase}
  \defaultfontfeatures[\rmfamily]{Ligatures=TeX,Scale=1}
\fi
\usepackage{lmodern}
\ifPDFTeX\else
  % xetex/luatex font selection
\fi
% Use upquote if available, for straight quotes in verbatim environments
\IfFileExists{upquote.sty}{\usepackage{upquote}}{}
\IfFileExists{microtype.sty}{% use microtype if available
  \usepackage[]{microtype}
  \UseMicrotypeSet[protrusion]{basicmath} % disable protrusion for tt fonts
}{}
\makeatletter
\@ifundefined{KOMAClassName}{% if non-KOMA class
  \IfFileExists{parskip.sty}{%
    \usepackage{parskip}
  }{% else
    \setlength{\parindent}{0pt}
    \setlength{\parskip}{6pt plus 2pt minus 1pt}}
}{% if KOMA class
  \KOMAoptions{parskip=half}}
\makeatother
\usepackage{xcolor}
\usepackage[margin=1in]{geometry}
\usepackage{color}
\usepackage{fancyvrb}
\newcommand{\VerbBar}{|}
\newcommand{\VERB}{\Verb[commandchars=\\\{\}]}
\DefineVerbatimEnvironment{Highlighting}{Verbatim}{commandchars=\\\{\}}
% Add ',fontsize=\small' for more characters per line
\newenvironment{Shaded}{}{}
\newcommand{\AlertTok}[1]{\textcolor[rgb]{1.00,0.00,0.00}{\textbf{#1}}}
\newcommand{\AnnotationTok}[1]{\textcolor[rgb]{0.38,0.63,0.69}{\textbf{\textit{#1}}}}
\newcommand{\AttributeTok}[1]{\textcolor[rgb]{0.49,0.56,0.16}{#1}}
\newcommand{\BaseNTok}[1]{\textcolor[rgb]{0.25,0.63,0.44}{#1}}
\newcommand{\BuiltInTok}[1]{\textcolor[rgb]{0.00,0.50,0.00}{#1}}
\newcommand{\CharTok}[1]{\textcolor[rgb]{0.25,0.44,0.63}{#1}}
\newcommand{\CommentTok}[1]{\textcolor[rgb]{0.38,0.63,0.69}{\textit{#1}}}
\newcommand{\CommentVarTok}[1]{\textcolor[rgb]{0.38,0.63,0.69}{\textbf{\textit{#1}}}}
\newcommand{\ConstantTok}[1]{\textcolor[rgb]{0.53,0.00,0.00}{#1}}
\newcommand{\ControlFlowTok}[1]{\textcolor[rgb]{0.00,0.44,0.13}{\textbf{#1}}}
\newcommand{\DataTypeTok}[1]{\textcolor[rgb]{0.56,0.13,0.00}{#1}}
\newcommand{\DecValTok}[1]{\textcolor[rgb]{0.25,0.63,0.44}{#1}}
\newcommand{\DocumentationTok}[1]{\textcolor[rgb]{0.73,0.13,0.13}{\textit{#1}}}
\newcommand{\ErrorTok}[1]{\textcolor[rgb]{1.00,0.00,0.00}{\textbf{#1}}}
\newcommand{\ExtensionTok}[1]{#1}
\newcommand{\FloatTok}[1]{\textcolor[rgb]{0.25,0.63,0.44}{#1}}
\newcommand{\FunctionTok}[1]{\textcolor[rgb]{0.02,0.16,0.49}{#1}}
\newcommand{\ImportTok}[1]{\textcolor[rgb]{0.00,0.50,0.00}{\textbf{#1}}}
\newcommand{\InformationTok}[1]{\textcolor[rgb]{0.38,0.63,0.69}{\textbf{\textit{#1}}}}
\newcommand{\KeywordTok}[1]{\textcolor[rgb]{0.00,0.44,0.13}{\textbf{#1}}}
\newcommand{\NormalTok}[1]{#1}
\newcommand{\OperatorTok}[1]{\textcolor[rgb]{0.40,0.40,0.40}{#1}}
\newcommand{\OtherTok}[1]{\textcolor[rgb]{0.00,0.44,0.13}{#1}}
\newcommand{\PreprocessorTok}[1]{\textcolor[rgb]{0.74,0.48,0.00}{#1}}
\newcommand{\RegionMarkerTok}[1]{#1}
\newcommand{\SpecialCharTok}[1]{\textcolor[rgb]{0.25,0.44,0.63}{#1}}
\newcommand{\SpecialStringTok}[1]{\textcolor[rgb]{0.73,0.40,0.53}{#1}}
\newcommand{\StringTok}[1]{\textcolor[rgb]{0.25,0.44,0.63}{#1}}
\newcommand{\VariableTok}[1]{\textcolor[rgb]{0.10,0.09,0.49}{#1}}
\newcommand{\VerbatimStringTok}[1]{\textcolor[rgb]{0.25,0.44,0.63}{#1}}
\newcommand{\WarningTok}[1]{\textcolor[rgb]{0.38,0.63,0.69}{\textbf{\textit{#1}}}}
\usepackage{longtable,booktabs,array}
\usepackage{calc} % for calculating minipage widths
% Correct order of tables after \paragraph or \subparagraph
\usepackage{etoolbox}
\makeatletter
\patchcmd\longtable{\par}{\if@noskipsec\mbox{}\fi\par}{}{}
\makeatother
% Allow footnotes in longtable head/foot
\IfFileExists{footnotehyper.sty}{\usepackage{footnotehyper}}{\usepackage{footnote}}
\makesavenoteenv{longtable}
\setlength{\emergencystretch}{3em} % prevent overfull lines
\providecommand{\tightlist}{%
  \setlength{\itemsep}{0pt}\setlength{\parskip}{0pt}}
\setcounter{secnumdepth}{5}
\ifLuaTeX
  \usepackage{selnolig}  % disable illegal ligatures
\fi
\usepackage{bookmark}
\IfFileExists{xurl.sty}{\usepackage{xurl}}{} % add URL line breaks if available
\urlstyle{same}
\hypersetup{
  colorlinks=true,
  linkcolor={blue},
  filecolor={Maroon},
  citecolor={Blue},
  urlcolor={blue},
  pdfcreator={LaTeX via pandoc}}

\author{}
\date{}

\begin{document}

{
\hypersetup{linkcolor=}
\setcounter{tocdepth}{3}
}
\section{SEQ BirdDex Cyberdeck --- Project Task
Sheet}\label{seq-birddex-cyberdeck-project-task-sheet}

\textbf{Project:} SEQ BirdDex\\
\textbf{Purpose:} Build an offline, battery-powered bird-identification
cyberdeck for South-East Queensland and adjacent inland/coastal
regions.\\
\textbf{Target trip:} September Lamington field use.\\
\textbf{Current status:} Concept / planning.\\
\textbf{Primary design constraint:} Field-useful by September, not
research-perfect.

\begin{center}\rule{0.5\linewidth}{0.5pt}\end{center}

\subsection{1. Executive summary}\label{executive-summary}

The SEQ BirdDex is a handheld or shoulder-carried cyberdeck that can
identify birds from field images without internet access. It should
accept photos from a Sony A7R V, take its own close-range photos, crop
the bird subject automatically, classify the likely bird species,
display a Pokédex-style species card, and log the sighting with time,
GPS, local weather, image path, crop path, prediction confidence, and
user confirmation.

The correct September target is not a perfect all-Australia classifier.
The correct target is a robust offline system that works well for common
and likely species in the defined region of interest, provides top-5
suggestions, handles uncertainty honestly, and creates a useful field
log.

The recommended v1 stack is:

\begin{Shaded}
\begin{Highlighting}[]
\NormalTok{Raspberry Pi 5 16GB}
\NormalTok{+ Raspberry Pi AI HAT+ 26 TOPS or 13 TOPS}
\NormalTok{+ NVMe SSD}
\NormalTok{+ daylight{-}readable touchscreen}
\NormalTok{+ local Wi{-}Fi access point}
\NormalTok{+ FTP ingest from Sony A7R V}
\NormalTok{+ local camera for close{-}range photos}
\NormalTok{+ GPS + temperature/humidity/pressure sensor}
\NormalTok{+ detector + classifier + geotemporal re{-}ranker}
\NormalTok{+ SQLite sighting database}
\end{Highlighting}
\end{Shaded}

The device should store the trained model, labels, species cards,
regional priors, reference thumbnails, and logs. It should not store the
full training image dataset unless a retrieval-based fallback is added
later.

\begin{center}\rule{0.5\linewidth}{0.5pt}\end{center}

\subsection{2. Region of interest}\label{region-of-interest}

The initial region of interest is the user-drawn SEQ / southern
Queensland region:

\begin{Shaded}
\begin{Highlighting}[]
\NormalTok{North: Bundaberg / Agnes Water}
\NormalTok{East: Fraser Coast, Sunshine Coast, Brisbane, Gold Coast, Tweed coast}
\NormalTok{South: Tweed Heads / northern NSW edge / possible Byron{-}Lismore buffer}
\NormalTok{West: Goondiwindi / Miles / Tara / Dalby corridor}
\NormalTok{Central: Gympie, Maryborough, Toowoomba{-}region hinterland, Scenic Rim, Lamington}
\end{Highlighting}
\end{Shaded}

The uploaded ROI image is a concept sketch, not yet a machine-readable
boundary. The first technical task is to convert the red boundary into a
GeoJSON polygon that can be used for API queries and occurrence
filtering.

\subsubsection{2.1 Approximate ROI anchor points for initial GeoJSON
draft}\label{approximate-roi-anchor-points-for-initial-geojson-draft}

These are seed anchors only. They must be checked manually in QGIS or
another map tool before data collection.

\begin{longtable}[]{@{}
  >{\raggedright\arraybackslash}p{(\columnwidth - 4\tabcolsep) * \real{0.3000}}
  >{\raggedleft\arraybackslash}p{(\columnwidth - 4\tabcolsep) * \real{0.4000}}
  >{\raggedright\arraybackslash}p{(\columnwidth - 4\tabcolsep) * \real{0.3000}}@{}}
\toprule\noalign{}
\begin{minipage}[b]{\linewidth}\raggedright
Anchor
\end{minipage} & \begin{minipage}[b]{\linewidth}\raggedleft
Approximate location
\end{minipage} & \begin{minipage}[b]{\linewidth}\raggedright
Purpose
\end{minipage} \\
\midrule\noalign{}
\endhead
\bottomrule\noalign{}
\endlastfoot
A & Agnes Water / Burnett coast & Northern coastal cap \\
B & Offshore Fraser / Hervey Bay buffer & Coastal record inclusion \\
C & Sunshine Coast / Noosa / Caloundra & Coastal SEQ corridor \\
D & Gold Coast / Tweed Heads & Southern coastal edge \\
E & Northern NSW buffer near Byron/Lismore & Optional overflow/vagrant
buffer \\
F & Scenic Rim / Lamington & Core September field zone \\
G & Goondiwindi & South-western inland edge \\
H & Miles / Tara / Dalby corridor & Western dryland/inland edge \\
I & North Burnett / Biggenden hinterland & North-western return
boundary \\
\end{longtable}

\subsubsection{2.2 ROI outputs}\label{roi-outputs}

Deliverables:

\begin{Shaded}
\begin{Highlighting}[]
\NormalTok{data/geo/seq\_birddex\_roi\_v0.geojson}
\NormalTok{data/geo/seq\_birddex\_roi\_v1\_reviewed.geojson}
\NormalTok{data/geo/seq\_birddex\_roi\_buffered\_25km.geojson}
\NormalTok{reports/roi\_definition.md}
\end{Highlighting}
\end{Shaded}

Acceptance criteria:

\begin{itemize}
\tightlist
\item
  The polygon can be loaded into QGIS.
\item
  The polygon can be used in Python via \texttt{shapely} /
  \texttt{geopandas}.
\item
  A 25 km buffer variant exists for vagrants and slightly out-of-bound
  observations.
\item
  The ROI can filter occurrence records from ALA, eBird-derived files,
  and iNaturalist exports.
\end{itemize}

\begin{center}\rule{0.5\linewidth}{0.5pt}\end{center}

\subsection{3. Product definition}\label{product-definition}

\subsubsection{3.1 User story}\label{user-story}

A user sees a bird in the field. They either photograph it with the Sony
A7R V or use the cyberdeck camera. The image arrives on the cyberdeck.
The cyberdeck detects the bird, crops it, predicts likely species,
displays the top candidates and a readable species card, and logs the
sighting offline. The user can confirm, reject, or mark the result as
unsure.

\subsubsection{3.2 Core functions}\label{core-functions}

\begin{longtable}[]{@{}llr@{}}
\toprule\noalign{}
ID & Function & Required for September? \\
\midrule\noalign{}
\endhead
\bottomrule\noalign{}
\endlastfoot
F-001 & Take local photos from cyberdeck camera & Yes \\
F-002 & Receive JPEGs from Sony A7R V & Yes \\
F-003 & Watch ingest folder for new images & Yes \\
F-004 & Detect bird subject in full-frame image & Yes \\
F-005 & Crop bird subject & Yes \\
F-006 & Classify cropped bird image & Yes \\
F-007 & Display top-5 predictions & Yes \\
F-008 & Display species card & Yes \\
F-009 & Log timestamp, GPS, weather, image, crop, prediction & Yes \\
F-010 & Let user confirm / reject / mark unsure & Yes \\
F-011 & Export sightings to CSV/JSON & Yes \\
F-012 & Run fully offline & Yes \\
F-013 & Audio bird ID & Optional v1.5 \\
F-014 & Map view of sightings & Optional \\
F-015 & Full Australasian species coverage & No, later \\
F-016 & RAW file classification & No, use JPEG for v1 \\
\end{longtable}

\subsubsection{3.3 Non-functional
requirements}\label{non-functional-requirements}

\begin{longtable}[]{@{}
  >{\raggedright\arraybackslash}p{(\columnwidth - 4\tabcolsep) * \real{0.3333}}
  >{\raggedright\arraybackslash}p{(\columnwidth - 4\tabcolsep) * \real{0.3333}}
  >{\raggedright\arraybackslash}p{(\columnwidth - 4\tabcolsep) * \real{0.3333}}@{}}
\toprule\noalign{}
\begin{minipage}[b]{\linewidth}\raggedright
ID
\end{minipage} & \begin{minipage}[b]{\linewidth}\raggedright
Requirement
\end{minipage} & \begin{minipage}[b]{\linewidth}\raggedright
Target
\end{minipage} \\
\midrule\noalign{}
\endhead
\bottomrule\noalign{}
\endlastfoot
NFR-001 & Offline operation & No internet required after setup \\
NFR-002 & Battery runtime & Minimum 4 h; stretch target 8 h \\
NFR-003 & Time to result & Target \textless{} 5 s per image; stretch
\textless{} 2 s \\
NFR-004 & Storage robustness & No database corruption after normal
shutdown \\
NFR-005 & Field UX & No terminal required during use \\
NFR-006 & Uncertainty handling & Must allow ``unknown/unsure'' state \\
NFR-007 & Image retention & Store original JPEG and crop \\
NFR-008 & Reproducibility & Model version and dataset manifest logged \\
NFR-009 & Maintainability & Repo must run from clean setup scripts \\
\end{longtable}

\begin{center}\rule{0.5\linewidth}{0.5pt}\end{center}

\subsection{4. System architecture}\label{system-architecture}

\subsubsection{4.1 High-level pipeline}\label{high-level-pipeline}

\begin{Shaded}
\begin{Highlighting}[]
\NormalTok{Sony A7R V / cyberdeck camera / manual upload}
\NormalTok{        |}
\NormalTok{        v}
\NormalTok{Image ingest service}
\NormalTok{        |}
\NormalTok{        v}
\NormalTok{Metadata extractor}
\NormalTok{  {-} EXIF timestamp}
\NormalTok{  {-} camera model}
\NormalTok{  {-} focal length if present}
\NormalTok{  {-} GPS if present}
\NormalTok{        |}
\NormalTok{        v}
\NormalTok{Bird detector}
\NormalTok{  {-} finds one or more bird boxes}
\NormalTok{        |}
\NormalTok{        v}
\NormalTok{Cropper}
\NormalTok{  {-} crops subject}
\NormalTok{  {-} pads crop safely}
\NormalTok{  {-} stores crop image}
\NormalTok{        |}
\NormalTok{        v}
\NormalTok{Species classifier}
\NormalTok{  {-} returns class probabilities}
\NormalTok{        |}
\NormalTok{        v}
\NormalTok{Geotemporal prior}
\NormalTok{  {-} ROI occurrence}
\NormalTok{  {-} month/season}
\NormalTok{  {-} habitat/climate}
\NormalTok{  {-} local checklist}
\NormalTok{        |}
\NormalTok{        v}
\NormalTok{Decision layer}
\NormalTok{  {-} top{-}5}
\NormalTok{  {-} confidence}
\NormalTok{  {-} unknown threshold}
\NormalTok{        |}
\NormalTok{        v}
\NormalTok{User interface + SQLite log}
\end{Highlighting}
\end{Shaded}

\subsubsection{4.2 Device service
architecture}\label{device-service-architecture}

\begin{Shaded}
\begin{Highlighting}[]
\NormalTok{birddex{-}ingest.service}
\NormalTok{  Watches FTP/upload/camera folders for new images.}

\NormalTok{birddex{-}infer.service}
\NormalTok{  Runs detection, cropping, classification, re{-}ranking, and logging.}

\NormalTok{birddex{-}ui.service}
\NormalTok{  Runs local touchscreen/web UI.}

\NormalTok{birddex{-}sensors.service}
\NormalTok{  Reads GPS, pressure, temperature, humidity, optional light sensor.}

\NormalTok{birddex{-}backup.service}
\NormalTok{  Exports sightings and model metadata to USB/NVMe backup.}
\end{Highlighting}
\end{Shaded}

\subsubsection{4.3 Ingest modes}\label{ingest-modes}

\begin{longtable}[]{@{}
  >{\raggedright\arraybackslash}p{(\columnwidth - 4\tabcolsep) * \real{0.3333}}
  >{\raggedright\arraybackslash}p{(\columnwidth - 4\tabcolsep) * \real{0.3333}}
  >{\raggedright\arraybackslash}p{(\columnwidth - 4\tabcolsep) * \real{0.3333}}@{}}
\toprule\noalign{}
\begin{minipage}[b]{\linewidth}\raggedright
Mode
\end{minipage} & \begin{minipage}[b]{\linewidth}\raggedright
Path
\end{minipage} & \begin{minipage}[b]{\linewidth}\raggedright
Notes
\end{minipage} \\
\midrule\noalign{}
\endhead
\bottomrule\noalign{}
\endlastfoot
Sony FTP JPEG & A7R V -\textgreater{} Pi Wi-Fi AP -\textgreater{} FTP
folder & Preferred field path \\
USB copy & Camera/card -\textgreater{} cyberdeck & Backup path \\
Pi camera & Local camera capture & Close-range or demo path \\
Phone upload & Phone browser -\textgreater{} local web UI & Optional \\
\end{longtable}

\subsubsection{4.4 RAW vs JPEG decision}\label{raw-vs-jpeg-decision}

Use JPEG for real-time recognition.

Rationale:

\begin{itemize}
\tightlist
\item
  RAW/ARW requires demosaicing, white balance, colour conversion, and
  resizing before inference.
\item
  RAW transfer is slower and wastes battery.
\item
  Sony A7R V can keep RAW on the camera card while sending a smaller
  JPEG to the cyberdeck.
\item
  For model training, RAW files can be archived later, but the field
  classifier should operate on JPEGs.
\end{itemize}

\begin{center}\rule{0.5\linewidth}{0.5pt}\end{center}

\subsection{5. Hardware recommendation}\label{hardware-recommendation}

\subsubsection{5.1 Recommended September
build}\label{recommended-september-build}

\begin{longtable}[]{@{}
  >{\raggedright\arraybackslash}p{(\columnwidth - 4\tabcolsep) * \real{0.3333}}
  >{\raggedright\arraybackslash}p{(\columnwidth - 4\tabcolsep) * \real{0.3333}}
  >{\raggedright\arraybackslash}p{(\columnwidth - 4\tabcolsep) * \real{0.3333}}@{}}
\toprule\noalign{}
\begin{minipage}[b]{\linewidth}\raggedright
Subsystem
\end{minipage} & \begin{minipage}[b]{\linewidth}\raggedright
Recommended part
\end{minipage} & \begin{minipage}[b]{\linewidth}\raggedright
Reason
\end{minipage} \\
\midrule\noalign{}
\endhead
\bottomrule\noalign{}
\endlastfoot
Compute & Raspberry Pi 5 16GB & Best ecosystem, camera support, GPIO,
documentation \\
AI acceleration & Raspberry Pi AI HAT+ 26 TOPS & Strong edge vision
acceleration with Hailo-8 \\
Storage & 256GB--1TB NVMe SSD & Fast logs/images/model storage \\
Screen & 5--7 inch touchscreen, high brightness preferred & Field UI \\
Local camera & Pi Camera Module 3 / HQ Camera / AI Camera & Close-range
capture and testing \\
External camera & Sony A7R V JPEG FTP transfer & High-quality bird
photos \\
GPS & u-blox USB/serial GPS & Location/time logging \\
Weather & BME280/BME680/SHT31 & Temperature/humidity/pressure \\
Power & USB-C PD battery bank or Li-ion pack with BMS & Field runtime \\
Controls & Physical buttons + touchscreen & Glove/sunlight-friendly
UI \\
Enclosure & 3D-printed cyberdeck shell & Custom ergonomics and
mounting \\
\end{longtable}

\subsubsection{5.2 Hardware alternatives}\label{hardware-alternatives}

\begin{longtable}[]{@{}
  >{\raggedright\arraybackslash}p{(\columnwidth - 6\tabcolsep) * \real{0.2500}}
  >{\raggedright\arraybackslash}p{(\columnwidth - 6\tabcolsep) * \real{0.2500}}
  >{\raggedright\arraybackslash}p{(\columnwidth - 6\tabcolsep) * \real{0.2500}}
  >{\raggedright\arraybackslash}p{(\columnwidth - 6\tabcolsep) * \real{0.2500}}@{}}
\toprule\noalign{}
\begin{minipage}[b]{\linewidth}\raggedright
Platform
\end{minipage} & \begin{minipage}[b]{\linewidth}\raggedright
Strengths
\end{minipage} & \begin{minipage}[b]{\linewidth}\raggedright
Weaknesses
\end{minipage} & \begin{minipage}[b]{\linewidth}\raggedright
Recommendation
\end{minipage} \\
\midrule\noalign{}
\endhead
\bottomrule\noalign{}
\endlastfoot
Pi 5 + AI HAT+ & Low software risk, strong ecosystem, Hailo
acceleration, good camera/GPIO support & Model export to Hailo can be
annoying & Best v1 platform \\
Jetson Orin Nano Super & Much stronger ML flexibility, CUDA/TensorRT, 67
TOPS class device & Higher power, more heat, less handheld-friendly &
Best high-performance dev platform \\
Orange Pi 5 Plus / RK3588 & Strong SBC CPU/NPU specs & NPU software
friction, weaker ecosystem & Avoid for deadline unless experimenting \\
Laptop/tablet backend & Very flexible & Not cyberdeck/off-grid elegant &
Useful for early dev only \\
Phone app only & Best finished bird ID experience & Not custom/offline
cyberdeck & Use as validation, not core build \\
\end{longtable}

\subsubsection{5.3 Power budget target}\label{power-budget-target}

Approximate v1 power planning values:

\begin{longtable}[]{@{}
  >{\raggedright\arraybackslash}p{(\columnwidth - 4\tabcolsep) * \real{0.3000}}
  >{\raggedleft\arraybackslash}p{(\columnwidth - 4\tabcolsep) * \real{0.4000}}
  >{\raggedright\arraybackslash}p{(\columnwidth - 4\tabcolsep) * \real{0.3000}}@{}}
\toprule\noalign{}
\begin{minipage}[b]{\linewidth}\raggedright
Component
\end{minipage} & \begin{minipage}[b]{\linewidth}\raggedleft
Expected draw class
\end{minipage} & \begin{minipage}[b]{\linewidth}\raggedright
Notes
\end{minipage} \\
\midrule\noalign{}
\endhead
\bottomrule\noalign{}
\endlastfoot
Pi 5 & Medium & Depends heavily on CPU load, screen, peripherals \\
AI HAT+ & Low to medium & Accelerator is efficient, but full system
still matters \\
Screen & Medium to high & Often the dominant field drain \\
NVMe & Low to medium & Spikes during writes \\
GPS/weather & Low & Negligible relative to screen/compute \\
Wi-Fi AP & Low to medium & Required for A7R ingest \\
Camera & Low to medium & Depends on module and use \\
\end{longtable}

Runtime design rules:

\begin{itemize}
\tightlist
\item
  Screen timeout after 30--60 s.
\item
  Classify on image arrival or button press, not continuous heavy
  inference.
\item
  Do not run large transformer inference continuously.
\item
  Use a graceful shutdown button.
\item
  Log battery voltage if using a custom pack.
\end{itemize}

\subsubsection{5.4 Hardware decision gate}\label{hardware-decision-gate}

Choose Pi 5 + AI HAT+ if:

\begin{itemize}
\tightlist
\item
  September field reliability matters more than peak ML speed.
\item
  You want Raspberry Pi camera/GPIO/support to work cleanly.
\item
  You are willing to tune/export models to Hailo or keep an ONNX/TFLite
  fallback.
\end{itemize}

Choose Jetson Orin Nano Super if:

\begin{itemize}
\tightlist
\item
  You want to run larger YOLO/ConvNeXt/ViT models locally.
\item
  You accept 7--25 W class power behaviour and extra thermal design.
\item
  You want CUDA/TensorRT more than Pi ecosystem convenience.
\end{itemize}

\begin{center}\rule{0.5\linewidth}{0.5pt}\end{center}

\subsection{6. Dataset plan}\label{dataset-plan}

\subsubsection{6.1 Dataset principle}\label{dataset-principle}

For v1, use supervised labelled data. Do not start with SimCLR as the
main model path.

Reason:

\begin{itemize}
\tightlist
\item
  A species classifier needs species labels.
\item
  Self-supervised learning can improve representations later, but it
  does not create species names by itself.
\item
  The fastest reliable path is transfer learning from pretrained image
  models, fine-tuned on labelled Australian/SEQ bird images.
\end{itemize}

\subsubsection{6.2 Dataset sources}\label{dataset-sources}

\begin{longtable}[]{@{}
  >{\raggedright\arraybackslash}p{(\columnwidth - 4\tabcolsep) * \real{0.3333}}
  >{\raggedright\arraybackslash}p{(\columnwidth - 4\tabcolsep) * \real{0.3333}}
  >{\raggedright\arraybackslash}p{(\columnwidth - 4\tabcolsep) * \real{0.3333}}@{}}
\toprule\noalign{}
\begin{minipage}[b]{\linewidth}\raggedright
Source
\end{minipage} & \begin{minipage}[b]{\linewidth}\raggedright
Use
\end{minipage} & \begin{minipage}[b]{\linewidth}\raggedright
Notes
\end{minipage} \\
\midrule\noalign{}
\endhead
\bottomrule\noalign{}
\endlastfoot
iNaturalist Australia & Labelled bird images, metadata, observations &
Track image license and attribution \\
Atlas of Living Australia & Occurrence data, taxonomy, species profiles,
spatial filtering & Good for ROI occurrence filtering \\
eBird / EBD / Status and Trends & Occurrence, seasonality, checklists,
relative abundance & Strong for geotemporal priors \\
BirdLife/Birdata & Australian bird context and occurrence records &
Useful for local validity checks \\
User A7R V photos & Domain adaptation and validation & Most valuable
after field trips \\
\end{longtable}

\subsubsection{6.3 Licensing requirements}\label{licensing-requirements}

Every downloaded image must store:

\begin{Shaded}
\begin{Highlighting}[]
\NormalTok{source}
\NormalTok{observation\_id}
\NormalTok{photo\_id}
\NormalTok{common\_name}
\NormalTok{scientific\_name}
\NormalTok{taxon\_id}
\NormalTok{license}
\NormalTok{creator/attribution}
\NormalTok{source\_url}
\NormalTok{date\_observed}
\NormalTok{latitude/longitude if available and non{-}sensitive}
\NormalTok{quality\_grade}
\end{Highlighting}
\end{Shaded}

Do not train on or redistribute images unless the license permits the
intended use. For personal experiments this is lower risk, but the
manifest should still preserve attribution and license metadata.

\subsubsection{6.4 Initial species scope}\label{initial-species-scope}

Use tiers.

\begin{longtable}[]{@{}
  >{\raggedright\arraybackslash}p{(\columnwidth - 4\tabcolsep) * \real{0.3333}}
  >{\raggedright\arraybackslash}p{(\columnwidth - 4\tabcolsep) * \real{0.3333}}
  >{\raggedright\arraybackslash}p{(\columnwidth - 4\tabcolsep) * \real{0.3333}}@{}}
\toprule\noalign{}
\begin{minipage}[b]{\linewidth}\raggedright
Tier
\end{minipage} & \begin{minipage}[b]{\linewidth}\raggedright
Scope
\end{minipage} & \begin{minipage}[b]{\linewidth}\raggedright
Target use
\end{minipage} \\
\midrule\noalign{}
\endhead
\bottomrule\noalign{}
\endlastfoot
Tier 1 & 100--200 common/likely ROI species & September MVP \\
Tier 2 & 300--500 broader QLD / northern NSW species & Post-MVP
expansion \\
Tier 3 & 700+ Australian regular species & Long-term goal \\
Tier 4 & Vagrants / rare records / subspecies / sex-age morphs &
Research-grade extension \\
\end{longtable}

\subsubsection{6.5 Dataset quality rules}\label{dataset-quality-rules}

\begin{itemize}
\tightlist
\item
  Split by observation/user/date where possible, not naive random image
  split.
\item
  Remove near-duplicates across train/validation/test.
\item
  Prefer images where bird is visible and labelled to species.
\item
  Keep difficult images for validation, not only pretty field-guide
  shots.
\item
  Track class imbalance.
\item
  Group visually similar species for confusion analysis.
\end{itemize}

\subsubsection{6.6 Data deliverables}\label{data-deliverables}

\begin{Shaded}
\begin{Highlighting}[]
\NormalTok{data/manifests/roi\_species\_candidates.parquet}
\NormalTok{data/manifests/species\_priority\_tiers.csv}
\NormalTok{data/manifests/images\_manifest.parquet}
\NormalTok{data/splits/train.csv}
\NormalTok{data/splits/val.csv}
\NormalTok{data/splits/test.csv}
\NormalTok{data/reports/license\_report.md}
\NormalTok{data/reports/class\_balance\_report.md}
\NormalTok{data/reports/duplicate\_audit.md}
\end{Highlighting}
\end{Shaded}

\begin{center}\rule{0.5\linewidth}{0.5pt}\end{center}

\subsection{7. Model plan}\label{model-plan}

\subsubsection{7.1 Recommended v1 model
architecture}\label{recommended-v1-model-architecture}

\begin{Shaded}
\begin{Highlighting}[]
\NormalTok{Detector:}
\NormalTok{  YOLO{-}style lightweight bird detector}
\NormalTok{  Purpose: find birds in raw photos}

\NormalTok{Classifier:}
\NormalTok{  EfficientNetV2{-}B0/B1 or MobileNetV3{-}Large}
\NormalTok{  Purpose: classify cropped bird image}

\NormalTok{Re{-}ranker:}
\NormalTok{  Visual probability + location prior + month prior + habitat prior}

\NormalTok{Decision layer:}
\NormalTok{  Top{-}5 display + confidence thresholds + unknown state}
\end{Highlighting}
\end{Shaded}

\subsubsection{7.2 Detector task}\label{detector-task}

The detector is not required to identify species. It only needs to find
bird-shaped objects.

Inputs:

\begin{Shaded}
\begin{Highlighting}[]
\NormalTok{Full{-}resolution or resized JPEG}
\end{Highlighting}
\end{Shaded}

Outputs:

\begin{Shaded}
\begin{Highlighting}[]
\NormalTok{bounding boxes}
\NormalTok{box confidence}
\NormalTok{crop image paths}
\end{Highlighting}
\end{Shaded}

Acceptance criteria:

\begin{itemize}
\tightlist
\item
  Detects birds in field images where the bird is small but visible.
\item
  Handles branches, sky, water, grass, and clutter.
\item
  Can return multiple birds.
\item
  Does not crop so tightly that tail/beak/wings are lost.
\end{itemize}

\subsubsection{7.3 Classifier task}\label{classifier-task}

Inputs:

\begin{Shaded}
\begin{Highlighting}[]
\NormalTok{Bird crop, usually 224–384 px input resolution}
\end{Highlighting}
\end{Shaded}

Outputs:

\begin{Shaded}
\begin{Highlighting}[]
\NormalTok{species probabilities}
\NormalTok{top{-}5 species}
\NormalTok{confidence / entropy / uncertainty}
\end{Highlighting}
\end{Shaded}

Candidate v1 backbones:

\begin{longtable}[]{@{}
  >{\raggedright\arraybackslash}p{(\columnwidth - 4\tabcolsep) * \real{0.3333}}
  >{\raggedright\arraybackslash}p{(\columnwidth - 4\tabcolsep) * \real{0.3333}}
  >{\raggedright\arraybackslash}p{(\columnwidth - 4\tabcolsep) * \real{0.3333}}@{}}
\toprule\noalign{}
\begin{minipage}[b]{\linewidth}\raggedright
Backbone
\end{minipage} & \begin{minipage}[b]{\linewidth}\raggedright
Strength
\end{minipage} & \begin{minipage}[b]{\linewidth}\raggedright
Weakness
\end{minipage} \\
\midrule\noalign{}
\endhead
\bottomrule\noalign{}
\endlastfoot
MobileNetV3-Large & Very efficient, mobile-friendly & May underperform
on fine-grained species \\
EfficientNetV2-B0/B1 & Good accuracy/speed compromise & Export and
quantisation need care \\
ConvNeXt-Tiny & Stronger visual features & Heavier for Pi \\
ViT/Swin & Strong with enough data/compute & Not ideal for Pi v1 \\
BioCLIP-style embedding & Useful for biological similarity and retrieval
& Heavier and more complex deployment \\
\end{longtable}

\subsubsection{7.4 Geotemporal re-ranking}\label{geotemporal-re-ranking}

The model should not rely on visual score alone. The system should
re-rank predictions using local probability.

Example scoring:

\begin{Shaded}
\begin{Highlighting}[]
\NormalTok{final\_score(species) =}
\NormalTok{  visual\_logit\_score}
\NormalTok{  + alpha * roi\_occurrence\_prior}
\NormalTok{  + beta  * month\_prior}
\NormalTok{  + gamma * habitat\_prior}
\NormalTok{  + delta * user\_history\_prior}
\end{Highlighting}
\end{Shaded}

Rules:

\begin{itemize}
\tightlist
\item
  Do not remove rare/vagrant species completely.
\item
  Penalise unlikely species rather than making them impossible.
\item
  Provide a ``vagrant/rare but possible'' warning when applicable.
\item
  Use a toggle for ``strict local mode'' versus ``wide/vagrant mode''.
\end{itemize}

\subsubsection{7.5 Unknown and uncertainty
logic}\label{unknown-and-uncertainty-logic}

The device must be allowed to say ``not sure.'' Suggested rules:

\begin{Shaded}
\begin{Highlighting}[]
\NormalTok{If max\_probability \textless{} threshold\_low:}
\NormalTok{  display Unknown / uncertain}

\NormalTok{If top1\_probability {-} top2\_probability \textless{} margin\_threshold:}
\NormalTok{  display Similar species warning}

\NormalTok{If visual score high but geotemporal prior low:}
\NormalTok{  display Possible rare/vagrant warning}

\NormalTok{If detector crop confidence low:}
\NormalTok{  display Crop uncertain warning}
\end{Highlighting}
\end{Shaded}

\subsubsection{7.6 Optional embedding
fallback}\label{optional-embedding-fallback}

A retrieval module may be added later.

\begin{Shaded}
\begin{Highlighting}[]
\NormalTok{Image crop {-}\textgreater{} embedding model {-}\textgreater{} nearest reference examples {-}\textgreater{} species vote}
\end{Highlighting}
\end{Shaded}

Device storage:

\begin{Shaded}
\begin{Highlighting}[]
\NormalTok{reference\_embeddings.faiss}
\NormalTok{embedding\_metadata.sqlite}
\NormalTok{reference\_thumbnails/}
\end{Highlighting}
\end{Shaded}

Use case:

\begin{itemize}
\tightlist
\item
  Low-confidence classifier output.
\item
  Showing similar reference birds.
\item
  Adding new species with fewer examples.
\end{itemize}

Do not make this required for September unless the classifier path is
already stable.

\begin{center}\rule{0.5\linewidth}{0.5pt}\end{center}

\subsection{8. Device data model}\label{device-data-model}

\subsubsection{8.1 SQLite schema}\label{sqlite-schema}

\begin{Shaded}
\begin{Highlighting}[]
\KeywordTok{CREATE} \KeywordTok{TABLE}\NormalTok{ model\_versions (}
    \KeywordTok{id} \DataTypeTok{INTEGER} \KeywordTok{PRIMARY} \KeywordTok{KEY}\NormalTok{,}
\NormalTok{    model\_name TEXT }\KeywordTok{NOT} \KeywordTok{NULL}\NormalTok{,}
\NormalTok{    model\_type TEXT }\KeywordTok{NOT} \KeywordTok{NULL}\NormalTok{,}
\NormalTok{    version TEXT }\KeywordTok{NOT} \KeywordTok{NULL}\NormalTok{,}
\NormalTok{    file\_path TEXT }\KeywordTok{NOT} \KeywordTok{NULL}\NormalTok{,}
\NormalTok{    created\_at TEXT,}
\NormalTok{    training\_manifest\_hash TEXT,}
\NormalTok{    notes TEXT}
\NormalTok{);}

\KeywordTok{CREATE} \KeywordTok{TABLE}\NormalTok{ species (}
\NormalTok{    species\_id TEXT }\KeywordTok{PRIMARY} \KeywordTok{KEY}\NormalTok{,}
\NormalTok{    birddex\_id TEXT }\KeywordTok{UNIQUE} \KeywordTok{NOT} \KeywordTok{NULL}\NormalTok{,}
\NormalTok{    common\_name TEXT }\KeywordTok{NOT} \KeywordTok{NULL}\NormalTok{,}
\NormalTok{    scientific\_name TEXT }\KeywordTok{NOT} \KeywordTok{NULL}\NormalTok{,}
\NormalTok{    family TEXT,}
\NormalTok{    group\_name TEXT,}
\NormalTok{    size\_text TEXT,}
\NormalTok{    habitat\_text TEXT,}
\NormalTok{    climate\_text TEXT,}
\NormalTok{    distribution\_text TEXT,}
\NormalTok{    similar\_species\_json TEXT,}
\NormalTok{    field\_marks\_text TEXT,}
\NormalTok{    facts\_text TEXT,}
\NormalTok{    conservation\_status TEXT,}
\NormalTok{    reference\_image\_path TEXT}
\NormalTok{);}

\KeywordTok{CREATE} \KeywordTok{TABLE}\NormalTok{ observations (}
    \KeywordTok{id} \DataTypeTok{INTEGER} \KeywordTok{PRIMARY} \KeywordTok{KEY}\NormalTok{,}
\NormalTok{    timestamp\_local TEXT }\KeywordTok{NOT} \KeywordTok{NULL}\NormalTok{,}
\NormalTok{    timestamp\_utc TEXT,}
\NormalTok{    latitude }\DataTypeTok{REAL}\NormalTok{,}
\NormalTok{    longitude }\DataTypeTok{REAL}\NormalTok{,}
\NormalTok{    altitude\_m }\DataTypeTok{REAL}\NormalTok{,}
\NormalTok{    temperature\_c }\DataTypeTok{REAL}\NormalTok{,}
\NormalTok{    humidity\_pct }\DataTypeTok{REAL}\NormalTok{,}
\NormalTok{    pressure\_hpa }\DataTypeTok{REAL}\NormalTok{,}
\NormalTok{    light\_lux }\DataTypeTok{REAL}\NormalTok{,}
    \KeywordTok{source}\NormalTok{ TEXT,}
\NormalTok{    original\_image\_path TEXT,}
\NormalTok{    crop\_image\_path TEXT,}
\NormalTok{    detector\_confidence }\DataTypeTok{REAL}\NormalTok{,}
\NormalTok{    top1\_species\_id TEXT,}
\NormalTok{    top1\_common\_name TEXT,}
\NormalTok{    top1\_confidence }\DataTypeTok{REAL}\NormalTok{,}
\NormalTok{    top5\_json TEXT,}
\NormalTok{    model\_version\_id }\DataTypeTok{INTEGER}\NormalTok{,}
\NormalTok{    user\_confirmed\_species\_id TEXT,}
\NormalTok{    user\_verdict TEXT,}
\NormalTok{    notes TEXT,}
    \KeywordTok{FOREIGN} \KeywordTok{KEY}\NormalTok{(model\_version\_id) }\KeywordTok{REFERENCES}\NormalTok{ model\_versions(}\KeywordTok{id}\NormalTok{)}
\NormalTok{);}

\KeywordTok{CREATE} \KeywordTok{TABLE}\NormalTok{ field\_sessions (}
    \KeywordTok{id} \DataTypeTok{INTEGER} \KeywordTok{PRIMARY} \KeywordTok{KEY}\NormalTok{,}
\NormalTok{    session\_name TEXT,}
\NormalTok{    start\_time TEXT,}
\NormalTok{    end\_time TEXT,}
\NormalTok{    location\_label TEXT,}
\NormalTok{    notes TEXT}
\NormalTok{);}
\end{Highlighting}
\end{Shaded}

\subsubsection{8.2 File layout on
cyberdeck}\label{file-layout-on-cyberdeck}

\begin{Shaded}
\begin{Highlighting}[]
\NormalTok{/opt/birddex/}
\NormalTok{  app/}
\NormalTok{  models/}
\NormalTok{    detector.hef}
\NormalTok{    classifier.hef}
\NormalTok{    classifier\_fallback.onnx}
\NormalTok{    label\_map.json}
\NormalTok{  data/}
\NormalTok{    birddex.sqlite}
\NormalTok{    priors.sqlite}
\NormalTok{    taxonomy.json}
\NormalTok{  media/}
\NormalTok{    original/}
\NormalTok{    crops/}
\NormalTok{    reference/}
\NormalTok{  logs/}
\NormalTok{  exports/}
\NormalTok{  config/}
\NormalTok{    device.yaml}
\NormalTok{    thresholds.yaml}
\NormalTok{    roi.geojson}
\end{Highlighting}
\end{Shaded}

\begin{center}\rule{0.5\linewidth}{0.5pt}\end{center}

\subsection{9. User interface
requirements}\label{user-interface-requirements}

\subsubsection{9.1 Main screens}\label{main-screens}

\begin{longtable}[]{@{}
  >{\raggedright\arraybackslash}p{(\columnwidth - 2\tabcolsep) * \real{0.5000}}
  >{\raggedright\arraybackslash}p{(\columnwidth - 2\tabcolsep) * \real{0.5000}}@{}}
\toprule\noalign{}
\begin{minipage}[b]{\linewidth}\raggedright
Screen
\end{minipage} & \begin{minipage}[b]{\linewidth}\raggedright
Required content
\end{minipage} \\
\midrule\noalign{}
\endhead
\bottomrule\noalign{}
\endlastfoot
Home & Camera status, battery, GPS lock, latest sighting \\
Incoming photo & Original image, detected crop(s), processing status \\
Prediction & Top-5 list, confidence, warnings \\
Species card & Name, BirdDex ID, reference image, facts, habitat,
climate \\
Confirm sighting & Confirm / wrong / unsure / add note \\
Logbook & Recent observations, search/filter \\
Export & CSV/JSON export, backup status \\
Settings & Model version, thresholds, ROI mode, storage \\
\end{longtable}

\subsubsection{9.2 Prediction display
format}\label{prediction-display-format}

Example:

\begin{Shaded}
\begin{Highlighting}[]
\NormalTok{Likely ID}
\NormalTok{BDX{-}0042 — Regent Bowerbird}
\NormalTok{Confidence: 0.72}
\NormalTok{Local likelihood: High}

\NormalTok{Other candidates:}
\NormalTok{2. Satin Bowerbird — 0.14}
\NormalTok{3. Green Catbird — 0.06}
\NormalTok{4. Golden Whistler — 0.03}
\NormalTok{5. Eastern Whipbird — 0.02}

\NormalTok{Warnings:}
\NormalTok{{-} Good crop}
\NormalTok{{-} Similar species possible}
\NormalTok{{-} Location/month plausible}

\NormalTok{Actions:}
\NormalTok{[Confirm] [Wrong] [Unsure] [Show similar]}
\end{Highlighting}
\end{Shaded}

\subsubsection{9.3 Physical controls}\label{physical-controls}

Recommended controls:

\begin{longtable}[]{@{}ll@{}}
\toprule\noalign{}
Button & Function \\
\midrule\noalign{}
\endhead
\bottomrule\noalign{}
\endlastfoot
Capture & Take local photo \\
Process latest & Run inference on latest received image \\
Confirm & Mark prediction correct \\
Unsure & Store for later review \\
Back & Return to previous screen \\
Power & Safe shutdown / wake \\
\end{longtable}

Touchscreen is useful, but physical buttons are more reliable in bright
outdoor conditions.

\begin{center}\rule{0.5\linewidth}{0.5pt}\end{center}

\subsection{10. Training and evaluation
plan}\label{training-and-evaluation-plan}

\subsubsection{10.1 Desktop training flow}\label{desktop-training-flow}

\begin{Shaded}
\begin{Highlighting}[]
\NormalTok{1. Build ROI species list.}
\NormalTok{2. Download labelled images with license metadata.}
\NormalTok{3. Clean dataset.}
\NormalTok{4. Build train/val/test splits.}
\NormalTok{5. Train baseline classifier.}
\NormalTok{6. Evaluate baseline.}
\NormalTok{7. Add detector/cropper.}
\NormalTok{8. Evaluate raw{-}image pipeline.}
\NormalTok{9. Calibrate confidence.}
\NormalTok{10. Export model to deployment format.}
\NormalTok{11. Run Pi{-}side inference tests.}
\end{Highlighting}
\end{Shaded}

\subsubsection{10.2 Metrics}\label{metrics}

\begin{longtable}[]{@{}
  >{\raggedright\arraybackslash}p{(\columnwidth - 4\tabcolsep) * \real{0.3333}}
  >{\raggedright\arraybackslash}p{(\columnwidth - 4\tabcolsep) * \real{0.3333}}
  >{\raggedright\arraybackslash}p{(\columnwidth - 4\tabcolsep) * \real{0.3333}}@{}}
\toprule\noalign{}
\begin{minipage}[b]{\linewidth}\raggedright
Metric
\end{minipage} & \begin{minipage}[b]{\linewidth}\raggedright
Purpose
\end{minipage} & \begin{minipage}[b]{\linewidth}\raggedright
Target for MVP
\end{minipage} \\
\midrule\noalign{}
\endhead
\bottomrule\noalign{}
\endlastfoot
Top-1 accuracy & Direct correctness & Useful but not sole metric \\
Top-5 accuracy & Field shortlist quality & High priority \\
Per-class accuracy & Find weak species & Required \\
Confusion matrix & Similar species clusters & Required \\
Expected calibration error & Confidence honesty & Required if
possible \\
Unknown rejection precision & Avoid false certainty & Required \\
Detector recall & Bird found in image & High priority \\
Time per image & Field usability & \textless{} 5 s target \\
Battery runtime & Field usability & \textgreater{} 4 h minimum \\
\end{longtable}

\subsubsection{10.3 Test sets}\label{test-sets}

Create multiple test sets:

\begin{Shaded}
\begin{Highlighting}[]
\NormalTok{clean\_test:}
\NormalTok{  curated labelled images from known sources}

\NormalTok{field\_like\_test:}
\NormalTok{  messy images, branches, distance, bad lighting}

\NormalTok{own\_camera\_test:}
\NormalTok{  A7R V photos from real field use}

\NormalTok{negative\_test:}
\NormalTok{  images without birds, signs, trees, dogs, insects, mammals}

\NormalTok{similar\_species\_test:}
\NormalTok{  honeyeaters, robins, thornbills, bowerbirds, fairywrens, raptors}
\end{Highlighting}
\end{Shaded}

\subsubsection{10.4 Acceptance criteria for model
v1}\label{acceptance-criteria-for-model-v1}

The model is acceptable for September if:

\begin{itemize}
\tightlist
\item
  Top-5 results are useful for common ROI birds.
\item
  The device can say ``unsure'' when confidence is low.
\item
  Similar species are grouped sensibly.
\item
  The cropper works on real field images.
\item
  Inference runs offline on the chosen device.
\item
  Every prediction logs model version and confidence.
\end{itemize}

\begin{center}\rule{0.5\linewidth}{0.5pt}\end{center}

\subsection{11. Development roadmap}\label{development-roadmap}

\subsubsection{Phase 0 --- Project skeleton and
ROI}\label{phase-0-project-skeleton-and-roi}

Tasks:

\begin{itemize}
\tightlist
\item[$\square$]
  Create Git repository.
\item[$\square$]
  Add project README.
\item[$\square$]
  Add \texttt{data/}, \texttt{src/}, \texttt{device/}, \texttt{models/},
  \texttt{reports/}, \texttt{tests/} layout.
\item[$\square$]
  Convert ROI sketch into GeoJSON.
\item[$\square$]
  Add ROI visualisation script.
\item[$\square$]
  Query initial species occurrence list.
\item[$\square$]
  Produce Tier 1 species list.
\end{itemize}

Deliverables:

\begin{Shaded}
\begin{Highlighting}[]
\NormalTok{README.md}
\NormalTok{roi.geojson}
\NormalTok{reports/roi\_definition.md}
\NormalTok{data/manifests/roi\_species\_candidates.csv}
\end{Highlighting}
\end{Shaded}

\subsubsection{Phase 1 --- Dataset
builder}\label{phase-1-dataset-builder}

Tasks:

\begin{itemize}
\tightlist
\item[$\square$]
  Implement iNaturalist/ALA metadata collector.
\item[$\square$]
  Implement image downloader with license filter.
\item[$\square$]
  Store source/license/attribution metadata.
\item[$\square$]
  Build duplicate detection workflow.
\item[$\square$]
  Build species balance report.
\item[$\square$]
  Build train/val/test splits.
\end{itemize}

Deliverables:

\begin{Shaded}
\begin{Highlighting}[]
\NormalTok{src/datasets/build\_species\_list.py}
\NormalTok{src/datasets/download\_images.py}
\NormalTok{src/datasets/build\_splits.py}
\NormalTok{data/manifests/images\_manifest.parquet}
\NormalTok{data/reports/license\_report.md}
\NormalTok{data/reports/class\_balance\_report.md}
\end{Highlighting}
\end{Shaded}

\subsubsection{Phase 2 --- Baseline
classifier}\label{phase-2-baseline-classifier}

Tasks:

\begin{itemize}
\tightlist
\item[$\square$]
  Train MobileNetV3 or EfficientNetV2 baseline.
\item[$\square$]
  Add augmentation pipeline.
\item[$\square$]
  Add class weighting or sampler for imbalance.
\item[$\square$]
  Save model checkpoints.
\item[$\square$]
  Generate evaluation report.
\item[$\square$]
  Export model to ONNX/TFLite.
\end{itemize}

Deliverables:

\begin{Shaded}
\begin{Highlighting}[]
\NormalTok{models/classifier\_v0.pt}
\NormalTok{models/classifier\_v0.onnx}
\NormalTok{reports/classifier\_v0\_eval.md}
\NormalTok{reports/confusion\_matrix\_v0.png}
\end{Highlighting}
\end{Shaded}

\subsubsection{Phase 3 --- Detector and
cropper}\label{phase-3-detector-and-cropper}

Tasks:

\begin{itemize}
\tightlist
\item[$\square$]
  Select detector model.
\item[$\square$]
  Test detector on field-like bird photos.
\item[$\square$]
  Implement crop padding and aspect-ratio handling.
\item[$\square$]
  Store crop previews.
\item[$\square$]
  Evaluate detection failures.
\end{itemize}

Deliverables:

\begin{Shaded}
\begin{Highlighting}[]
\NormalTok{models/detector\_v0.onnx}
\NormalTok{src/inference/detect\_and\_crop.py}
\NormalTok{reports/crop\_quality\_report.md}
\end{Highlighting}
\end{Shaded}

\subsubsection{Phase 4 --- End-to-end desktop
inference}\label{phase-4-end-to-end-desktop-inference}

Tasks:

\begin{itemize}
\tightlist
\item[$\square$]
  Build \texttt{predict\_image.py}.
\item[$\square$]
  Input full image, output crop + top-5.
\item[$\square$]
  Add geotemporal re-ranking stub.
\item[$\square$]
  Add JSON result export.
\item[$\square$]
  Add batch inference mode for test folders.
\end{itemize}

Deliverables:

\begin{Shaded}
\begin{Highlighting}[]
\NormalTok{src/inference/predict\_image.py}
\NormalTok{src/inference/rerank.py}
\NormalTok{reports/end\_to\_end\_eval.md}
\end{Highlighting}
\end{Shaded}

\subsubsection{Phase 5 --- Cyberdeck
runtime}\label{phase-5-cyberdeck-runtime}

Tasks:

\begin{itemize}
\tightlist
\item[$\square$]
  Install OS.
\item[$\square$]
  Configure NVMe storage.
\item[$\square$]
  Configure local Wi-Fi AP.
\item[$\square$]
  Configure FTP server.
\item[$\square$]
  Configure systemd services.
\item[$\square$]
  Add image folder watcher.
\item[$\square$]
  Run inference on new files.
\item[$\square$]
  Write results to SQLite.
\end{itemize}

Deliverables:

\begin{Shaded}
\begin{Highlighting}[]
\NormalTok{device/systemd/birddex{-}ingest.service}
\NormalTok{device/systemd/birddex{-}infer.service}
\NormalTok{device/config/device.yaml}
\NormalTok{device/scripts/install\_device.sh}
\end{Highlighting}
\end{Shaded}

\subsubsection{Phase 6 --- UI}\label{phase-6-ui}

Tasks:

\begin{itemize}
\tightlist
\item[$\square$]
  Build local web UI or touchscreen UI.
\item[$\square$]
  Display latest image and crop.
\item[$\square$]
  Display top-5 predictions.
\item[$\square$]
  Display species card.
\item[$\square$]
  Add confirm/wrong/unsure buttons.
\item[$\square$]
  Add logbook view.
\item[$\square$]
  Add export button.
\end{itemize}

Deliverables:

\begin{Shaded}
\begin{Highlighting}[]
\NormalTok{device/ui/}
\NormalTok{device/api/}
\NormalTok{reports/ui\_test\_report.md}
\end{Highlighting}
\end{Shaded}

\subsubsection{Phase 7 --- Sensors and
logging}\label{phase-7-sensors-and-logging}

Tasks:

\begin{itemize}
\tightlist
\item[$\square$]
  Integrate GPS.
\item[$\square$]
  Integrate weather sensor.
\item[$\square$]
  Add timestamp fallback via RTC.
\item[$\square$]
  Add battery status if available.
\item[$\square$]
  Log sensor snapshot per observation.
\end{itemize}

Deliverables:

\begin{Shaded}
\begin{Highlighting}[]
\NormalTok{src/device/sensors.py}
\NormalTok{reports/sensor\_test\_report.md}
\end{Highlighting}
\end{Shaded}

\subsubsection{Phase 8 --- Field
hardening}\label{phase-8-field-hardening}

Tasks:

\begin{itemize}
\tightlist
\item[$\square$]
  Battery runtime test.
\item[$\square$]
  Thermal test.
\item[$\square$]
  Screen readability test.
\item[$\square$]
  FTP transfer reliability test.
\item[$\square$]
  Sudden power-loss test using copied database only.
\item[$\square$]
  Graceful shutdown test.
\item[$\square$]
  Enclosure revision.
\end{itemize}

Deliverables:

\begin{Shaded}
\begin{Highlighting}[]
\NormalTok{reports/battery\_runtime.md}
\NormalTok{reports/thermal\_test.md}
\NormalTok{reports/field\_trial\_01.md}
\NormalTok{enclosure/stl/v1/}
\end{Highlighting}
\end{Shaded}

\begin{center}\rule{0.5\linewidth}{0.5pt}\end{center}

\subsection{12. September MVP
definition}\label{september-mvp-definition}

The MVP is complete when:

\begin{itemize}
\tightlist
\item[$\square$]
  Device boots into UI without terminal work.
\item[$\square$]
  Device creates its own Wi-Fi AP.
\item[$\square$]
  Sony A7R V can send JPEGs to the device.
\item[$\square$]
  New JPEGs are processed automatically.
\item[$\square$]
  The bird subject is cropped automatically.
\item[$\square$]
  The classifier returns a top-5 list offline.
\item[$\square$]
  The UI shows a species card.
\item[$\square$]
  The user can confirm/wrong/unsure the prediction.
\item[$\square$]
  GPS/weather/time/image/crop/prediction are logged.
\item[$\square$]
  Logs can be exported after the trip.
\item[$\square$]
  Battery runtime is at least 4 hours in realistic field use.
\end{itemize}

The MVP is not required to:

\begin{itemize}
\tightlist
\item[$\square$]
  Identify every Australian bird.
\item[$\square$]
  Run large transformer models.
\item[$\square$]
  Classify RAW files directly.
\item[$\square$]
  Work perfectly on tiny distant birds.
\item[$\square$]
  Upload to eBird/Birdata/iNaturalist automatically.
\item[$\square$]
  Be waterproof.
\end{itemize}

\begin{center}\rule{0.5\linewidth}{0.5pt}\end{center}

\subsection{13. Suggested repository
structure}\label{suggested-repository-structure}

\begin{Shaded}
\begin{Highlighting}[]
\NormalTok{seq{-}birddex/}
\NormalTok{  README.md}
\NormalTok{  pyproject.toml}
\NormalTok{  requirements.txt}
\NormalTok{  .gitignore}

\NormalTok{  data/}
\NormalTok{    geo/}
\NormalTok{    manifests/}
\NormalTok{    splits/}
\NormalTok{    external/}
\NormalTok{    interim/}
\NormalTok{    processed/}

\NormalTok{  src/}
\NormalTok{    birddex/}
\NormalTok{      datasets/}
\NormalTok{      training/}
\NormalTok{      inference/}
\NormalTok{      priors/}
\NormalTok{      export/}
\NormalTok{      device/}
\NormalTok{      ui/}
\NormalTok{      utils/}

\NormalTok{  models/}
\NormalTok{    detector/}
\NormalTok{    classifier/}
\NormalTok{    exported/}

\NormalTok{  device/}
\NormalTok{    systemd/}
\NormalTok{    config/}
\NormalTok{    scripts/}
\NormalTok{    api/}
\NormalTok{    ui/}

\NormalTok{  species\_cards/}
\NormalTok{    cards/}
\NormalTok{    reference\_images/}
\NormalTok{    species\_cards.sqlite}

\NormalTok{  reports/}
\NormalTok{    figures/}
\NormalTok{    eval/}
\NormalTok{    field\_tests/}

\NormalTok{  tests/}
\NormalTok{    test\_dataset\_manifest.py}
\NormalTok{    test\_inference\_pipeline.py}
\NormalTok{    test\_sqlite\_schema.py}
\NormalTok{    test\_reranker.py}
\end{Highlighting}
\end{Shaded}

\begin{center}\rule{0.5\linewidth}{0.5pt}\end{center}

\subsection{14. Risk register}\label{risk-register}

\begin{longtable}[]{@{}
  >{\raggedright\arraybackslash}p{(\columnwidth - 6\tabcolsep) * \real{0.2143}}
  >{\raggedleft\arraybackslash}p{(\columnwidth - 6\tabcolsep) * \real{0.2857}}
  >{\raggedleft\arraybackslash}p{(\columnwidth - 6\tabcolsep) * \real{0.2857}}
  >{\raggedright\arraybackslash}p{(\columnwidth - 6\tabcolsep) * \real{0.2143}}@{}}
\toprule\noalign{}
\begin{minipage}[b]{\linewidth}\raggedright
Risk
\end{minipage} & \begin{minipage}[b]{\linewidth}\raggedleft
Severity
\end{minipage} & \begin{minipage}[b]{\linewidth}\raggedleft
Probability
\end{minipage} & \begin{minipage}[b]{\linewidth}\raggedright
Mitigation
\end{minipage} \\
\midrule\noalign{}
\endhead
\bottomrule\noalign{}
\endlastfoot
Dataset is noisy or mislabelled & High & High & Manual audit, use
research-grade sources, filter low-quality observations \\
Similar species confusion & High & High & Top-5 UI, similar-species
cards, geotemporal priors \\
Bird too small in frame & High & Medium & Use A7R V JPEGs, detector,
crop inspection, confidence gating \\
Hailo export friction & Medium & Medium & Keep ONNX/TFLite fallback;
test export early \\
Battery runtime poor & Medium & Medium & Screen timeout,
classify-on-demand, large battery, low-power mode \\
UI too fiddly outdoors & High & Medium & Physical buttons, large fonts,
high contrast \\
Scope creep & High & High & Lock v1 to ROI + Tier 1 species \\
Legal/licensing issues & Medium & Medium & Store license metadata, avoid
redistribution of restricted images \\
Corrupt logs after power loss & Medium & Low/Medium & SQLite WAL mode,
graceful shutdown, periodic backups \\
Overconfident wrong IDs & High & Medium & Calibration, unknown
threshold, confidence display \\
\end{longtable}

\begin{center}\rule{0.5\linewidth}{0.5pt}\end{center}

\subsection{15. Immediate next actions}\label{immediate-next-actions}

Start with these in order:

\begin{enumerate}
\def\labelenumi{\arabic{enumi}.}
\tightlist
\item
  Convert the red ROI sketch into \texttt{roi.geojson}.
\item
  Generate an initial species candidate list for the ROI.
\item
  Select Tier 1 species for September.
\item
  Build a licensed image metadata manifest.
\item
  Train a small baseline classifier on desktop.
\item
  Test classifier on ugly real-world bird photos.
\item
  Prototype the Sony A7R V to Pi JPEG transfer path.
\item
  Build the Pi folder watcher and SQLite logger.
\item
  Add the simple UI.
\item
  Only then optimise model acceleration.
\end{enumerate}

\begin{center}\rule{0.5\linewidth}{0.5pt}\end{center}

\subsection{16. Source notes}\label{source-notes}

The following sources should be checked during implementation, because
hardware/API support can change:

\begin{itemize}
\tightlist
\item
  Raspberry Pi AI HAT+ documentation:
  https://www.raspberrypi.com/documentation/accessories/ai-hat-plus.html
\item
  NVIDIA Jetson Orin Nano Super Developer Kit:
  https://www.nvidia.com/en-au/autonomous-machines/embedded-systems/jetson-orin/nano-super-developer-kit/
\item
  Sony Camera Remote SDK:
  https://support.d-imaging.sony.co.jp/app/sdk/en/index.html
\item
  Sony camera/mobile compatibility support:
  https://www.sony.com.au/electronics/support/articles/00284526
\item
  Atlas of Living Australia API documentation: https://docs.ala.org.au/
\item
  eBird API documentation:
  https://documenter.getpostman.com/view/664302/S1ENwy59
\item
  eBird data products:
  https://science.ebird.org/en/use-ebird-data/download-ebird-data-products
\item
  iNaturalist Australia: https://inaturalist.ala.org.au/
\item
  BirdNET: https://birdnet.cornell.edu/
\item
  BirdNET-Pi: https://www.birdweather.com/birdnetpi
\item
  BioCLIP 2 model card: https://huggingface.co/imageomics/bioclip-2
\end{itemize}

\begin{center}\rule{0.5\linewidth}{0.5pt}\end{center}

\subsection{17. Final design stance}\label{final-design-stance}

For this project, the device should prioritise reliability over
theoretical peak model accuracy. A compact detector/classifier with good
crops, good uncertainty handling, strong local species priors, and a
useful logbook will outperform a large fragile model in real field use.

The first hard proof should be:

\begin{Shaded}
\begin{Highlighting}[]
\NormalTok{A real A7R V bird JPEG arrives on the cyberdeck,}
\NormalTok{the bird is cropped automatically,}
\NormalTok{the offline classifier returns a plausible top{-}5,}
\NormalTok{the UI shows a species card,}
\NormalTok{and the sighting is saved with GPS/weather/time metadata.}
\end{Highlighting}
\end{Shaded}

That is the milestone that proves the project is real.

\end{document}
